% Avsnitt 16.10, ur lectures/L16/appendix/b_exercises.md.

\section{Övningar}
\label{c16:sec:exercises}

\begin{aside}
Du bygger sändningsvägen, och \code{can_controller_tb} driver två noder genom ett komplett utbyte, så
den kan inte passera förrän \chapref{c17:ch} lägger till mottagningshalvan. \code{make build-project}
vet det och rapporterar den som \emph{överhoppad} i stället för att köra den och misslyckas:

\begin{console}
==> can_controller_tb (skipped: can_controller.vhd has no receive path yet; ...)
\end{console}

Så bygget förblir grönt det här kapitlet, och de här övningarna resonerar sig igenom sändningsvägen
och analyserar den i GHDL i stället. Testbänken blir domen om ett kapitel.
\end{aside}

\exercise{Följ en kort ram för hand}{Resonemang}
\label{c16:ex:shortframe}
Skriv, utifrån ramformatet i \secref{c10:sec:frameformat} och det här kapitlets fälttabell, ut hela
följden av tillstånd som \code{can_controller} passerar (i TX-rollen) för att sända en ram med ID
\code{0x001}, DLC \code{0} (inget datafält), från \code{STATE_IDLE} tillbaka till \code{STATE_IDLE}.
Notera för varje tillstånd ungefär hur många bitperioder det upptar.

\exercise{Varför tillståndet \code{STATE_LOAD_*} finns}{Resonemang}
\label{c16:ex:loadstate}
\xpart{a} Förklara med egna ord vad som skulle gå fel om ett stoppat fält sekvenserades med ett enda
tillstånd som gav \code{tx_shift_reg}:s \code{load} \emph{och} väntade på \code{done}, i stället för
ett separat \code{STATE_LOAD_*} på en cykel följt av skifttillståndet.

\xpart{b} Peka ut den specifika tajmingregeln från \chapref{c14:ch} (att \code{tx_shift_reg}:s
\code{load} presenterar direkt, och dess \code{done} som kommer ett skift extra) som gör uppdelningen
i två tillstånd nödvändig.

\exercise{Vad som når \code{crc15}, och vad som inte gör det}{Resonemang}
\label{c16:ex:crcfeed}
\xpart{a} Sortera för en sänd ram de här kategorierna av bitar i ”matas in i \code{crc15}” och
”utesluts”: SOF-biten, riktiga bitar ur ID, kontrollfält och data, instoppade stoppbitar, CRC-fältets
egna bitar, och den ostoppade svansen (CRC-avgränsare, ACK-plats, ACK-avgränsare, EOF). Säg för varje
utesluten kategori vad som utesluter den: antingen en term i
\code{crc_enable <= txsr_bit_valid and not txsr_stuff and role;}, eller det faktum att fältet aldrig
passerar genom \code{tx_shift_reg} över huvud taget.

\xpart{b} En av de kategorierna kan överraska dig: CRC-fältets egna bitar går tillbaka in i motorn.
Säg, med \chapref{c13:ch}:s egenskap att generering och kontroll är samma operation, vilket värde
registret håller efter att den sista CRC-biten matats in, och namnge den kontroll i \chapref{c17:ch}
som beror på precis det värdet. Förklara sedan varför sändningsvägen, tagen för sig, inte beror på
det: vad garanterar \code{STATE_START}:s \code{clear} redan om den andra ram en nod sänder?

\xpart{c} Anta att du i stället hade grindat bort \code{crc_enable} under \code{STATE_CRC_HI} och
\code{STATE_CRC_LO}. Skulle den första ram en nod sänder bära en korrekt CRC? Skulle den andra? Var
noga: att \code{crc15} nollställs är en del av svaret. Säg sedan om någon simulering i den här boken
skulle fånga ändringen, och vad det säger dig om var det återmatade CRC-fältet faktiskt gör nytta.

\exercise{Hur lång är en ram, exakt?}{Resonemang}
\label{c16:ex:framelen}
Att svara på det kräver att hela fälttabellen och stoppbitsregeln hålls i huvudet samtidigt, vilket är
själva poängen. Ta ramen \code{can_controller} sänder för ID \code{0x123}, DLC \code{3}, databytesen
\code{FF 00 00}.

\xpart{a} Skriv ut det stoppade området bit för bit, SOF till och med datafältets slut, med det här
kapitlets fälttabell. (Det ska bli 43 bitar: $1 + 11 + 1 + 2 + 4 + 24$.) Identifieraren är
\code{0x123}; kom ihåg att RTR, IDE och r0 alla är \code{0} och att DLC är \code{0011}.

\xpart{b} CRC-15 över de 43 bitarna är \code{0x0FEE}, alltså \code{000111111101110} (du kan bekräfta
det mot din \code{crc15} från \chapref{c13:ch} om du vill). Lägg till den, vilket ger hela det
stoppade området på 58 bitar, SOF till och med CRC.

\xpart{c} Tillämpa stoppbitsregeln över alla 58 bitarna och räkna de instoppade stoppbitarna. Håll
koll på gränsen mellan datafältets avslutande nollor och CRC:ns inledande nollor: följdräknaren
nollställs inte vid fältgränser.

\xpart{d} Lägg till den ostoppade svansen (CRC-avgränsare, ACK-plats, ACK-avgränsare, EOF) och ange
det totala antalet bitar som faktiskt drivs ut på ledningen. Hur många mikrosekunder är det vid
1~Mbit/s?

\xpart{e} Samma ram med datan \code{11 22 33} behöver bara en stoppbit i stället för antalet du fick
i \textbf{c)}. Förklara med en mening varför en rams varaktighet på en CAN-buss beror på
\emph{nyttolastens värde} och inte bara på dess längd, och vad det innebär för den som försöker räkna
ut värsta fallets bustajming.
